\begin{table}[H]
\centering
\small
\begin{tabular}{lrrrrrrr}
\toprule
$\lambda$ & $\lVert\nabla J^\lambda-\nabla J\rVert$ & $/\lambda^2$ & $\lVert\E[\widehat G]-\nabla J^\lambda\rVert$ & $p$ & Dispersion & $d\cdot\mathrm{MSE}$ & $\cos$ \\
\midrule
$0.05$ & $0.0030$ & $1.18$ & $0.0421$ & $1.000$ & $0.5926$ & $0.3500$ & $0.941$ \\
$0.1$ & $0.0115$ & $1.15$ & $0.0480$ & $1.000$ & $0.6151$ & $0.3763$ & $0.914$ \\
$0.2$ & $0.0433$ & $1.08$ & $0.0702$ & $0.996$ & $0.7234$ & $0.5224$ & $0.892$ \\
$0.4$ & $0.1562$ & $0.98$ & $0.1963$ & $0.123$ & $1.2290$ & $1.5232$ & $0.742$ \\
$0.8$ & $0.5899$ & $0.92$ & $0.9565$ & $0.000$ & $4.7778$ & $23.2859$ & $0.643$ \\
\bottomrule
\end{tabular}
\caption{Evaluation-only gradient diagnostics on the linear--quadratic benchmark, at the policies saved at the end of training and averaged over the five seeds, from 50 independent replications of the estimator at 8192 particles. Column two is the perturbation bias of the gradient, $\lVert\nabla_\theta J^\lambda(\theta)-\nabla_\theta J(\theta)\rVert$, computed from the closed-form gradient oracle; column three divides it by $\lambda^2$. Column four is the bias of the estimator with respect to the perturbed gradient it targets, and column five the $p$-value of the $\chi^2$ test that this bias is zero at the available replication count, so a large $p$ means the estimator is not measurably biased for $\nabla_\theta J^\lambda$. Column six is the dispersion $\lVert\operatorname{sd}(\widehat G)\rVert$ of the replications and column seven the mean-square error scaled by the parameter dimension $d$; it equals the sum of the squares of columns four and six, up to the factor $(R-1)/R$ carried by the unbiased dispersion estimate, and the identity is verified run by run to $2\%$, which is exactly that factor at $R=50$. The last column is the cosine alignment of the mean estimate with the oracle gradient.}
\label{tab:gradient-diagnostics-lq}
\end{table}
